% GENERATED FILE. Edit canonical lessons, then run npm run generate:books.
% canonical-source-hash: fnv1a64:f753d0ff0f7dd227
% canonical-lessons: MR-C205-mophat, MR-C205-adhunik, MR-C205-pracin, MR-C205-sijlela, MR-C205-piklela

\chapter{Free of charge, Modern, Ancient, Cooked, Ripe}
\label{ch:mr-free-of-charge-modern-ancient-cooked-ripe}
\hlchaptermodality{205}

\begin{chapteropening}
\textbf{By the end of this chapter:} \emph{I can say free of charge, modern, ancient, cooked and ripe.}

Complete the last lesson of chapter 205: I can say free of charge, modern, ancient, cooked and ripe.
\end{chapteropening}

\section[\texorpdfstring{\mr{मोफत} (mophat)}{मोफत (mophat)}]{\mr{मोफत} (mophat) — free of charge}

\label{lesson:MR-C205-mophat}



\begin{quote}
Before the new one: say the Marathi for fearful, then the Marathi for curious.
\end{quote}

\subsection*{You'll want to know: \mr{मोफत}}
\textbf{\mr{मोफत}} — \emph{mophat} — \textquotedblleft{}free of charge\textquotedblright{}.

\begin{grammarlens}[title={What you've built}]
One more word for this chapter.
\end{grammarlens}

\subsection*{Your turn}
\begin{itemize}
\raggedright
  \item \emph{Say it:} \emph{mophat}
  \item \emph{Say it:} \emph{mophat}, once more
  \item \emph{Say it:} say \emph{mophat}
  \item \emph{Recall:} say the Marathi for fearful, then the Marathi for curious, then say \emph{mophat} again
\end{itemize}

\begin{tcolorbox}[breakable,colback=teal!4,colframe=teal!35!black,title={Before you move on}]
What does \mr{मोफत} mean? (\textquotedblleft{}Free of charge\textquotedblright{}.) Say it once more.
\end{tcolorbox}
\section[\texorpdfstring{\mr{आधुनिक} (ādhunik)}{आधुनिक (ādhunik)}]{\mr{आधुनिक} (ādhunik) — modern}

\label{lesson:MR-C205-adhunik}



\begin{quote}
Before the new one: say the Marathi for curious, then the Marathi for free of charge.
\end{quote}

\subsection*{You'll want to know: \mr{आधुनिक}}
\textbf{\mr{आधुनिक}} — \emph{ādhunik} — \textquotedblleft{}modern\textquotedblright{}.

\begin{grammarlens}[title={What you've built}]
One more word for this chapter.
\end{grammarlens}

\subsection*{Your turn}
\begin{itemize}
\raggedright
  \item \emph{Say it:} \emph{ādhunik}
  \item \emph{Say it:} \emph{ādhunik}, once more
  \item \emph{Say it:} say \emph{ādhunik}
  \item \emph{Recall:} say the Marathi for curious, then the Marathi for free of charge, then say \emph{ādhunik} again
\end{itemize}

\begin{tcolorbox}[breakable,colback=teal!4,colframe=teal!35!black,title={Before you move on}]
What does \mr{आधुनिक} mean? (\textquotedblleft{}Modern\textquotedblright{}.) Say it once more.
\end{tcolorbox}
\section[\texorpdfstring{\mr{प्राचीन} (prācīn)}{प्राचीन (prācīn)}]{\mr{प्राचीन} (prācīn) — ancient}

\label{lesson:MR-C205-pracin}



\begin{quote}
Before the new one: say the Marathi for free of charge, then the Marathi for modern.
\end{quote}

\subsection*{You'll want to know: \mr{प्राचीन}}
\textbf{\mr{प्राचीन}} — \emph{prācīn} — \textquotedblleft{}ancient\textquotedblright{}.

\begin{grammarlens}[title={What you've built}]
One more word for this chapter.
\end{grammarlens}

\subsection*{Your turn}
\begin{itemize}
\raggedright
  \item \emph{Say it:} \emph{prācīn}
  \item \emph{Say it:} \emph{prācīn}, once more
  \item \emph{Say it:} say \emph{prācīn}
  \item \emph{Recall:} say the Marathi for free of charge, then the Marathi for modern, then say \emph{prācīn} again
\end{itemize}

\begin{tcolorbox}[breakable,colback=teal!4,colframe=teal!35!black,title={Before you move on}]
What does \mr{प्राचीन} mean? (\textquotedblleft{}Ancient\textquotedblright{}.) Say it once more.
\end{tcolorbox}
\section[\texorpdfstring{\mr{शिजलेला} (śijlelā)}{शिजलेला (śijlelā)}]{\mr{शिजलेला} (śijlelā) — cooked}

\label{lesson:MR-C205-sijlela}



\begin{quote}
Before the new one: say the Marathi for modern, then the Marathi for ancient.
\end{quote}

\subsection*{You'll want to know: \mr{शिजलेला}}
\textbf{\mr{शिजलेला}} — \emph{śijlelā} — \textquotedblleft{}cooked\textquotedblright{}.

\begin{grammarlens}[title={What you've built}]
One more word for this chapter.
\end{grammarlens}

\subsection*{Your turn}
\begin{itemize}
\raggedright
  \item \emph{Say it:} \emph{śijlelā}
  \item \emph{Say it:} \emph{śijlelā}, once more
  \item \emph{Say it:} say \emph{śijlelā}
  \item \emph{Recall:} say the Marathi for modern, then the Marathi for ancient, then say \emph{śijlelā} again
\end{itemize}

\begin{tcolorbox}[breakable,colback=teal!4,colframe=teal!35!black,title={Before you move on}]
What does \mr{शिजलेला} mean? (\textquotedblleft{}Cooked\textquotedblright{}.) Say it once more.
\end{tcolorbox}
\section[\texorpdfstring{\mr{पिकलेला} (piklelā)}{पिकलेला (piklelā)}]{\mr{पिकलेला} (piklelā) — ripe}

\label{lesson:MR-C205-piklela}



\begin{quote}
Before the new one: say the Marathi for ancient, then the Marathi for cooked.
\end{quote}

\subsection*{You'll want to know: \mr{पिकलेला}}
\textbf{\mr{पिकलेला}} — \emph{piklelā} — \textquotedblleft{}ripe\textquotedblright{}.

\begin{grammarlens}[title={What you've built}]
One more word for this chapter.
\end{grammarlens}

\subsection*{Your turn}
\begin{itemize}
\raggedright
  \item \emph{Say it:} \emph{piklelā}
  \item \emph{Say it:} \emph{piklelā}, once more
  \item \emph{Say it:} say \emph{piklelā}
  \item \emph{Recall:} say the Marathi for ancient, then the Marathi for cooked, then say \emph{piklelā} again
\end{itemize}

\begin{tcolorbox}[breakable,colback=teal!4,colframe=teal!35!black,title={Before you move on}]
What does \mr{पिकलेला} mean? (\textquotedblleft{}Ripe\textquotedblright{}.) Say it once more.
\end{tcolorbox}
